  \item \field{Innovation Details and Method}

    \placeholder{Explain how the innovation works. This section should make the concept concrete enough for reviewers to understand the design, technology, workflow, and prototype plan.}

    \textbf{6.1 Design Approach}
    \begin{itemize}
      \item Inclusive Design / Universal Design principles applied in the project.
      \item Accessibility Design decisions that support people with disabilities and other user groups.
      \item Key user needs, constraints, and assumptions.
    \end{itemize}

    \textbf{6.2 Technology Used}
    \begin{itemize}
      \item Hardware and software components.
      \item AI, IoT, mobile application, sensor, cloud, or other technical systems used.
      \item Development tools, platforms, datasets, or models.
    \end{itemize}

    \textbf{6.3 System Workflow}

    \placeholder{Add diagrams or supporting figures where useful.}
    \begin{itemize}
      \item System diagram.
      \item Flowchart.
      \item System architecture.
      \item Prototype images.
      \item User journey.
    \end{itemize}

    \textbf{6.4 Innovation Strengths}

    \placeholder{Explain what makes the project different from existing solutions.}
    \begin{itemize}
      \item Ease of use.
      \item Affordability.
      \item Support for multiple disability groups.
      \item AI-assisted analysis or personalization.
      \item Potential to increase employment or income opportunities.
    \end{itemize}

    \placeholder{Add tables, charts, prototype screenshots, or experimental results here if available.}
